% Generated draft: requires mathematical and source review.
\documentclass[11pt]{article}
\usepackage[margin=25mm]{geometry}
\usepackage{fontspec}
\setmainfont{Latin Modern Roman}
\usepackage{amsmath,amssymb,mathtools,unicode-math}
\setmathfont{Latin Modern Math}
\usepackage{enumitem,needspace,xurl,hyperref}
\setlength{\parindent}{0pt}
\setlength{\parskip}{5pt}
\setlength{\emergencystretch}{3em}
\newcommand{\N}{\mathbb{N}}
\newcommand{\Z}{\mathbb{Z}}
\newcommand{\Q}{\mathbb{Q}}
\newcommand{\R}{\mathbb{R}}
\newcommand{\C}{\mathbb{C}}
\newcommand{\F}{\mathbb{F}}
\newcommand{\A}{\mathbb{A}}
\newcommand{\PP}{\mathbb{P}}
\newcommand{\E}{\mathbb{E}}
\newcommand{\eps}{\varepsilon}
\newcommand{\dd}{\,\mathrm{d}}
\begin{document}
\section*{UnsolvedMath ID: 30005363}
\textbf{Source problem code:} catalogue entry
\subsection{Definitions and mathematical statement}
\textbf{Unrecoverable mathematical statement.} The supplied record for permanent UnsolvedMath ID \(30005363\), source code \texttt{catalogue entry}, contains a table-of-contents fragment beginning with the isolated word ``conjecture'' and a page reference to 160, followed by other workshop contributions. It supplies no identifiable conjecture, definitions, hypotheses, domains, boundary conditions, quantifiers, or conclusion. Consequently, no precise mathematical statement can be formulated from this evidence.

The supplied \texttt{statement\_status} is \texttt{unrecoverable}, and no clean statement is supplied. The supplied verification note reports that the official report identifies the extraction as a table-of-contents fragment; this is a recorded verification claim, not an independent verification here. The registry's ``open'' label and the listed contribution titles do not establish an identifiable open mathematical problem. In particular, neither the dangling word ``conjecture'' nor any subsequent contribution can safely be selected as the intended target.
\subsection{Short English statement}
The record preserves a workshop contents fragment rather than an identifiable mathematical problem. A statement-only draft cannot recover the intended conjecture or its assumptions from the supplied evidence. The supplied ``open'' label remains a curation label requiring independent verification.
\subsection{Sources}
\begin{itemize}
\item Local UnsolvedMath registry: supplied complete record for permanent ID \(30005363\), source code \texttt{catalogue entry}, titled \emph{Workshop Contents Listing for Low-Dimensional Topology}, in the Oberwolfach Reports Open Problems collection. The supplied proposed year is 2023; no problem proposer is identified.
\item Supplied report citation: \emph{Low-dimensional Topology} (2023), Oberwolfach Reports, pp.\ 151--213, DOI 10.4171/owr/2023/3. Supplied source URL and literature-source link labelled ``original Oberwolfach report'': \url{https://doi.org/10.4171/owr/2023/3}. An authoritative report citation is supplied, but no identifiable mathematical problem statement is supplied from it.
\item Literature assessment dated exactly \texttt{2026-08-22}: the supplied curation labels the record ``open'' and assesses the conjecture fragment as containing no identifiable mathematical assertion whose solved or false status could be verified. This is supplied dated curation, not an independently verified literature assessment. The background explicitly requests independent verification of time-sensitive status and cited claims.
\item The supplied verification note says that the extraction is a table-of-contents fragment, that no separate clean statement is promoted, and that the original record remains publicly visible and retained for provenance. These are supplied source-verification qualifications.
\item Attribution retained from the adjacent contents entries, without assigning any of them authorship of the unidentified conjecture: Sarah Blackwell, joint with Robion Kirby, Michael Klug, Vincent Longo, and Benjamin Ruppik, \emph{How not to study low-dimensional topology}, p.\ 163; Nir Lazarovich, joint with Yoav Moriah and Tali Pinsky, \emph{Highly twisted diagrams}, p.\ 166; Adele Jackson, \emph{Triangulations of Seifert fibered spaces}, p.\ 169; Alexander He, joint with Ben Burton, \emph{Selectively searching the Pachner graph}, p.\ 170; Paula Truöl, \emph{Strongly quasipositive knots are concordant to infinitely many strongly quasipositive knots}, p.\ 172.
\item Further adjacent contents attributions: Marissa Loving, joint with Elizabeth Field, Heejoung Kim, Autumn Kent, and Christopher Leininger, \emph{End-periodic homeomorphisms and volumes of mapping tori}, p.\ 174; Andrew Yarmola, joint with Tommaso Cremaschi, Jake Intrater, and José Andrés Rodriguez-Migueles, \emph{On canonical knots in tangent bundles}, p.\ 178; Ivan Dynnikov, \emph{Comparing Legendrian knots: a general algorithm and practical solutions}, p.\ 180; Tali Pinsky, joint with Bronek Wajnryb and Eiko Kin, \emph{Orbits of homeomorphisms of T 2 and fibrations of the Whitehead link}, p.\ 182; David Futer, joint with David Gabai and Andrew Yarmola, \emph{In search of the Margulis constant}, p.\ 183; José Pedro Quintanilha, joint with Stefan Friedl, Lars Munser, and Yuri Santos-Rego, \emph{Canonical decompositions and algorithmic recognition of spatial graphs}, p.\ 187. The extracted spelling ``T 2'' is retained without supplying a mathematical interpretation.
\end{itemize}
\end{document}
